\part{Begrijpen}

% ─────────────────────────────────────────────────────────────────
\chapter{Hemelcoördinaten}
\label{ch:coords}

\section{De hemelbol}

Vanaf de grond gezien gedraagt de sterrenhemel zich als een bol die van oost
naar west draait om een as door de hemelpolen. De sterren staan op zeer
uiteenlopende afstanden, maar hun plaats op die denkbeeldige bol volstaat om
ze te vinden en te volgen.

De Aarde draait in \textbf{23~u 56~min 4~s} eenmaal om haar as --- de
\emph{sterrendag}, korter dan de 24 uur van de zonnedag. Dát is het tempo
waarop een montering moet draaien, en niet dat van de klok.

\section{Rechte klimming}

De \textbf{rechte klimming} (RK) is het hemelse equivalent van de lengtegraad.
Zij wordt geteld in \textbf{uren} en niet in graden, van 0~u tot 24~u, vanaf
het lentepunt.

\begin{formula}
\[
1\,\text{u} = 15\degre \qquad 1\,\text{min} = 15' \qquad 1\,\text{s} = 15''
\]
Zo is RK = 02~u 44~min 08~s gelijk aan $2\times15\degre + 44\times0{,}25\degre +
8\times\frac{15}{3600}\degre = 41{,}03\degre$.
\end{formula}

\begin{info}
RK is de kritieke coördinaat voor het volgen: een fout van één tijdseconde is
aan de hemelequator 15 boogseconden waard. Bij 1,6\arcsec/pixel is dat ongeveer
tien pixels streep.
\end{info}

\section{Declinatie}

De \textbf{declinatie} (DEC) is het equivalent van de breedtegraad: de
hoekafstand ten noorden of ten zuiden van de hemelequator, van $-90\degre$ tot
$+90\degre$. Op een correct opgestelde montering beweegt de declinatieas niet
tijdens het volgen.

\section{Uurhoek}

De \textbf{uurhoek} meet hoe ver een object van de plaatselijke meridiaan
staat. Hij groeit onophoudelijk met de draaiing van de Aarde.

\begin{formula}
\[
\text{uurhoek} = \text{plaatselijke sterrentijd} - \text{RK}
\]
Een RK-tolerantie laat zich uitdrukken in boogseconden of in tijdseconden. De
omrekening hangt af van de declinatie:
\[
\text{tijdseconden} = \frac{\text{boogseconden}}{15 \cos(\text{DEC})}
\]
\end{formula}

\begin{tip}
Dicht bij de hemelpool verplaatst een RK-fout de ster nauwelijks op de sensor:
de cosinus van de declinatie dempt haar. Circumpolaire objecten zijn dus
toegeeflijker voor volgfouten in RK --- maar niet voor die in declinatie.
\end{tip}

% ─────────────────────────────────────────────────────────────────
\chapter{Equatoriale monteringen}
\label{ch:mounts}

\section{Beginsel}

Een \textbf{equatoriale montering} heft de draaiing van de Aarde op door de
buis te draaien om een as evenwijdig aan die van de Aarde zelf: de poolas.

\begin{itemize}
  \item \textbf{Poolas (RK)}: draait op sterrensnelheid.
  \item \textbf{Declinatieas}: staat stil tijdens het volgen, in theorie.
\end{itemize}

\begin{formula}
\[
\omega = \frac{1\,296\,000''}{86\,164{,}09\,\text{s}} \approx
15{,}041\,\text{boogseconden per seconde}
\]
\end{formula}

Elke afwijking van dat tempo, en elke scheefstand van de poolas, wordt een
drift die de opname vastlegt.

\section{10Micron-monteringen}

De 10Micron-monteringen (GM1000 HPS en hoger) dragen \textbf{absolute
encoders}: zij weten waar zij staan, zelfs na een stroomonderbreking, waar een
relatieve encoder dat kwijtraakt. Daar komen een wormwielaandrijving, correctie
van de periodieke fout en vooral een intern \textbf{aanwijsmodel} bij, opgebouwd
uit enkele tientallen sterren.

Dat model maakt \textbf{ongegidst} fotograferen mogelijk: de montering draait
niet louter op sterrensnelheid, zij past gecorrigeerde snelheden toe die
doorbuiging, refractie en de resterende poolscheefstand opvangen.

\begin{warning}
Een rechtstreeks gevolg, dat vaak verkeerd begrepen wordt: op een montering met
model \textbf{zit de poolopstellingsfout ín het model} en wordt zij al door de
volgsnelheden gecompenseerd. Een resterende declinatiedrift vraagt dus niet om
de azimutschroef, maar om \textbf{meer modelpunten in de buurt van het doel}.
\end{warning}

Deze monteringen spreken een \textbf{uitgebreid LX200}: naast de standaard
commando's geven zij ruwe asposities, toegepaste snelheden, de staat van het
model en omgevingsgegevens door. MountMonitor gebruikt die uitbreidingen waar
ze beschikbaar zijn.

% ─────────────────────────────────────────────────────────────────
\chapter{Wat het volgen verslechtert}
\label{ch:errors}

Geen enkele montering volgt volmaakt. De oorzaken stapelen zich op, en elk
laat een eigen handtekening in de registratie achter.

\section{Periodieke fout}

Een wormwiel is nooit volmaakt rond en nooit volmaakt gecentreerd. Zijn afwijking
herhaalt zich bij elke omwenteling en levert een trilling met vaste periode op
--- meestal enkele minuten. Dat is precies wat een Fourier-analyse eruit kan
halen (hoofdstuk~\ref{ch:fft}).

\section{Poolopstelling}

Wijst de poolas niet precies naar de hemelpool, dan drift de ster langzaam in
declinatie, en des te sneller naarmate men verder van de meridiaan staat. Op
een montering zonder model is het antwoord mechanisch. Op een montering mét
model: zie de waarschuwing in het vorige hoofdstuk.

\section{Atmosferische refractie}

De dampkring tilt objecten des te meer op naarmate zij lager staan: zo'n $34'$
aan de horizon, $1'$ op $45\degre$ hoogte. De refractie verandert dus in de
loop van de nacht, en haar drift wordt gemakkelijk voor een scheefstand
aangezien.

\section{Mechanische verstoring en trillingen}

Wind, een trekkende kabel, bewegende ondergrond, een meridiaanommekeer: even
zovele eenmalige schokken, geen van alle periodiek. Zij tonen zich als
losstaande pieken, en juist daaraan zijn ze te herkennen.

\section{Aangestuurde bewegingen}

Dit zijn in het geheel geen fouten. Tussen twee opnames verplaatst de sequencer
de montering met opzet --- \textbf{dither} om de ruis te decorreleren,
\textbf{hercentrering} wanneer het doel is weggegleden --- en de montering
\emph{blijft} op haar nieuwe plaats. Ze verwarren met volgfouten is de duurste
vergissing die men bij het lezen van een registratie kan maken, en
hoofdstuk~\ref{ch:stats} is eraan gewijd.

% ─────────────────────────────────────────────────────────────────
\chapter{De cijfers lezen}
\label{ch:stats}

\begin{warning}
Dit is het belangrijkste hoofdstuk van de handleiding. Een volgregistratie
wordt met ontstellend gemak verkeerd gelezen, en het cijfer dat er dan uit komt
is een factor enkele honderden mis.
\end{warning}

\section{Het signaal is een trap}

Dit is wat een nacht ongegidst fotograferen werkelijk vastlegt: de montering
houdt haar positie op een fractie van een boogseconde, dan verplaatst de
sequencer haar met een dither, en zij houdt haar nieuwe positie even goed. Het
spoor is geen ruizige lijn om een gemiddelde: het is een \textbf{trap}.

Een standaardafwijking over de hele nacht meet dan de hoogte van de treden, en
niet de kwaliteit van het volgen. En het verwijderen van een rechte lijn ---
wat een klassieke ontdrifting doet --- verwijdert geen trap.

\begin{info}
Bij een werkelijke sessie van achtenhalf uur gaf de ruwe gecombineerde RMS
\textbf{86,3\arcsec}, een vernietigend oordeel. De montering hield in
werkelijkheid elke positie op \textbf{0,18\arcsec}. Tussen beide cijfers zit
een factor 480.
\end{info}

\section{De jitter}

De \textbf{jitter} is de spreiding die de montering vertoont \emph{terwijl} zij
een positie vasthoudt. Het is het enige cijfer dat een ster op de sensor
uitsmeert, en het cijfer waarop MountMonitor zijn oordeel baseert.

Hij wordt tussen twee herpositioneringen gemeten, nooit om een rechte lijn
heen. Het programma knipt de registratie bij de sprongen, laat twee meetpunten
na elke sprong weg --- de beweging zelf is geen volgen --- en neemt de mediaan
van de spreidingen per trede.

\begin{tip}
Deze werkwijze heeft een eigenschap die vermelding verdient: een reeks die bij
elke opname dithert, komt er niet slechter uit dan een reeks die dat nooit
doet. Het cijfer beschrijft de montering, niet de sequencer.
\end{tip}

\section{De aangestuurde bewegingen}

Zij worden apart geteld en ingedeeld:

\begin{center}
\rowcolors{2}{rowalt}{white}
\begin{tabularx}{\textwidth}{lX}
\toprule
\textbf{Klasse} & \textbf{Wat het is} \\
\midrule
Dither & Kleine verplaatsing tussen twee opnames, om de ruis te decorreleren. \\
Hercentrering & Ruimere verplaatsing, die het doel terugbrengt waar het hoort. \\
Niet gestabiliseerd & De spreiding die erop volgt blijft meerdere malen de
gewone waarde. Juist die verdient aandacht. \\
\bottomrule
\end{tabularx}
\end{center}

\begin{warning}
Een drempeloverschrijding is geen beweging. Een hercentrering van 70\arcsec{}
heeft meerdere meetpunten nodig om te voltooien en overschrijdt de drempel bij
elk daarvan: één nacht las zo 98 overschrijdingen voor 53 werkelijke
bewegingen, met een mediane amplitude die tien keer te klein was. MountMonitor
groepeert overschrijdingen die minder dan tien seconden uiteenliggen.
\end{warning}

\section{De langzame drift}

Hij wordt per uur en per opname gegeven. Bij een ongegidste montering heft elke
dither hem op: wat telt is niet de drift over de nacht, maar de verplaatsing
die hij \emph{tijdens één opname} veroorzaakt.

\begin{formula}
\[
\text{verplaatsing per opname} = \text{drift (\arcsec/u)} \times
\frac{\text{opnameduur (s)}}{3600}
\]
Een drift van $-16{,}8\arcsec$/u verplaatst de ster dus over $0{,}84\arcsec$
bij een opname van 180~s.
\end{formula}

\section{De lopende standaardafwijking}

Live berekend over een venster van 60 tot 900 seconden, dient zij om te zien
\emph{wanneer} er iets is gebeurd, niet om de nacht te beoordelen. Na een
doelwisseling, een onderbreking in de registratie of een herpositionering
overspant haar venster de beweging en beschrijft zij díe, niet het volgen: die
waarden blijven buiten het rapport.

\section{Het oordeel}

Het berust op de gecombineerde jitter, met drempels die bij het type montering
passen. Een ongegidste 10Micron wordt op ruimere drempels beoordeeld dan een
gegidste montering, want het zijn niet dezelfde ordes van grootte.

\begin{center}
\rowcolors{2}{rowalt}{white}
\begin{tabularx}{\textwidth}{lXX}
\toprule
\textbf{Oordeel} & \textbf{Gegidste montering} & \textbf{Ongegidst, precisie} \\
\midrule
Uitstekend & $< 0{,}5\arcsec$ & $< 2{,}0\arcsec$ \\
Goed       & $< 1{,}5\arcsec$ & $< 4{,}0\arcsec$ \\
Redelijk   & $< 3{,}0\arcsec$ & $< 8{,}0\arcsec$ \\
\bottomrule
\end{tabularx}
\end{center}

\begin{tip}
Betrek het oordeel altijd op uw bemonstering. Bij 1,6\arcsec/pixel laat een
jitter van 0,5\arcsec{} de ster een derde pixel trillen: onzichtbaar op de
uiteindelijke opname.
\end{tip}

% ─────────────────────────────────────────────────────────────────
\chapter{De Fourier-analyse}
\label{ch:fft}

\section{Waartoe zij dient}

Een Fourier-transformatie ontleedt een signaal in een som van sinussen. Op een
volgregistratie brengt zij naar voren wat zich \emph{herhaalt} --- en een
periodieke wormfout herhaalt zich per definitie.

MountMonitor zoekt in het bereik van \textbf{2 tot 15 minuten}, waarin de
wormperiodes van gangbare monteringen vallen.

\section{Twee voorzorgen}

\begin{itemize}
  \item De transformatie werkt op \textbf{ontdrifte} afwijkingen. Een lineaire
  drift is niet periodiek, maar over een eindig venster levert zij een grote
  laagfrequente component op die zou verbergen waarnaar gezocht wordt.
  \item De \textbf{grens van Nyquist} eist minstens twee metingen per periode.
  Bij 2~Hz zijn periodes onder de seconde buiten bereik --- geen bezwaar, geen
  enkele mechanische fout huist daar.
\end{itemize}

\section{Een spectrum lezen}

Een scherpe piek bij een stabiele periode wijst op een periodieke fout. Een vlak
spectrum betekent dat er geen aantoonbare is --- het geval van de
10Micron-monteringen na correctie. Brede, zwervende pieken zijn ruis, meestal
wind.

% ═══════════════════════════════════════════════════════════════════
\part{Gebruiken}

% ─────────────────────────────────────────────────────────────────
\chapter{Installatie}
\label{ch:install}

\section{Ondersteunde systemen}

De uitgegeven pakketten dragen hun eigen Python en hun eigen Qt mee: verder
hoeft er niets geïnstalleerd te worden. Wat zij wél vragen is een voldoende
recent systeem, omdat de Qt-binaries erin dat vragen.

\begin{center}
\rowcolors{2}{rowalt}{white}
\begin{tabularx}{\textwidth}{lXl}
\toprule
\textbf{Systeem} & \textbf{Minimaal} & \textbf{Architectuur} \\
\midrule
Windows & Windows 10 (1809) & x64 \\
macOS & macOS 11 Big Sur & Apple Silicon \emph{en} Intel \\
Linux & glibc 2.28 --- Debian 11, Ubuntu 20.04, RHEL 8 & x86\_64 \\
Linux ARM & glibc 2.39 --- Ubuntu 24.04 & arm64 \\
\bottomrule
\end{tabularx}
\end{center}

\begin{warning}
Kies het macOS-pakket dat bij uw machine past: \texttt{-macos-arm64.dmg} voor
een Mac met Apple Silicon, \texttt{-macos-x86\_64.dmg} voor een Intel-Mac. De
meegeleverde Python is architectuurgebonden, ook al is Qt zelf universeel: het
verkeerde pakket start in het geheel niet.
\end{warning}

\section{Het programma plaatsen}

Onder Windows regelt het installatieprogramma alles. Onder macOS sleept u de
toepassing naar de map Programma's. Onder Linux installeert
\texttt{installer.sh} haar in het gebruikersprofiel, zonder \texttt{sudo} en
zonder pakketbeheerder.

Updates worden daarna bij elke start stil gecontroleerd en aangeboden zodra er
zijn.

\section{Een eerste proef zonder montering}

Met de simulatiemodus loopt u alles door zonder apparatuur:

\begin{lstlisting}[language=bash]
MountMonitor --sim-all
\end{lstlisting}

% ─────────────────────────────────────────────────────────────────
\chapter{Instellingen}
\label{ch:config}

\section{Verbinding}

\begin{center}
\rowcolors{2}{rowalt}{white}
\begin{tabularx}{\textwidth}{lX}
\toprule
\textbf{Protocol} & \textbf{Wanneer te gebruiken} \\
\midrule
LX200 TCP/IP & Een 10Micron op het netwerk. Standaard poort 3492. \\
LX200 serieel & RS-232-verbinding, af fabriek 9600 8N1 bij 10Micron. \\
ASCOM & Alleen Windows, via het stuurprogramma van de fabrikant. \\
\bottomrule
\end{tabularx}
\end{center}

\section{De waarnemingslocatie}

\begin{warning}
Vul \textbf{breedtegraad, lengtegraad en hoogte} in bij de voorkeuren. Eerst
wordt de montering bevraagd (\texttt{:Gt\#} en \texttt{:Gg\#}), maar veel
opstellingen hebben haar nooit hun locatie doorgegeven en dan antwoordt zij
niets. Zonder locatie kunnen de efemeriden van de nacht niet worden berekend.

De lengtegraad voert u \textbf{positief naar het oosten} in. Het
LX200-protocol telt andersom; de omrekening gebeurt voor u zodra de montering
antwoordt.
\end{warning}

\section{Toleranties en lopend venster}

De tolerantie is een alarmdrempel die u zelf bepaalt. Een eenvoudige regel:
neem uw bemonstering in boogseconden per pixel. Een tolerantie gelijk aan de
bemonstering is streng, het dubbele is comfortabel.

Het venster van de lopende standaardafwijking stelt u in tussen 60 en 900
seconden. Zestig reageert snel, negenhonderd vlakt een lange sessie af.

\section{De logger gereedzetten}

Drie instellingen sturen de opnameknop aan. Alle drie staan \textbf{standaard
aan}.

\textbf{Registratie starten zodra de montering begint te volgen.} De knop zet
de logger dan \emph{gereed} in plaats van hem te starten: er wordt niets
geschreven zolang de montering niet volgt. Een montering die om twaalf uur
's middags wordt aangesloten voor een nacht die om 19~u begint, registreert
niet langer zeven uur niets. Een tweede klik schakelt uit.

\textbf{Onderbreken zodra de montering stopt met volgen}, met een respijttijd
van twee minuten. Die respijttijd bestaat met reden: een onderbreking kan van
voorbijgaande aard zijn --- een meridiaanommekeer, een autofocus, een zwenking
tussen twee doelen --- en afbreken bij het eerste meetpunt zou bij elke dither
een gat slaan.

\textbf{De nacht beëindigen bij zonsopkomst.} Een nacht is één geheel:
parkeert de montering om 02~u en gaat zij om 03~u verder, dan wordt de
registratie onderbroken en hervat \emph{in hetzelfde bestand}. Alleen het
daglicht sluit de sessie af en zet het rapport in gang --- en alleen wanneer de
zon tijdens de sessie werkelijk onder de horizon is geweest, zodat een proef
overdag zichzelf niet afsluit.

% ─────────────────────────────────────────────────────────────────
\chapter{Een nacht waarnemen}
\label{ch:usage}

\section{Het verloop}

\begin{enumerate}
  \item De montering verbinden en nagaan of de status op volgen springt.
  \item De logger gereedzetten --- of starten, als gereedzetten uitstaat.
  \item Laten lopen. Het rapport schrijft zichzelf aan het eind.
\end{enumerate}

\section{Waarop live te letten}

\begin{itemize}
  \item \textbf{De lopende standaardafwijking}, eerder dan de momentane
  afwijking: zij zegt of de montering haar positie vasthoudt.
  \item \textbf{De tolerantiemeldingen}, die een overschrijding aangeven.
  \item \textbf{De marge tot het dauwpunt}, als er een omgevingssensor is
  aangesloten.
\end{itemize}

\begin{tip}
Een losstaande piek is vrijwel nooit de montering: het is een windvlaag, een
kabel, een voetstap te dicht bij de zuil. Wat telt is het niveau \emph{tussen}
de pieken.
\end{tip}

\section{De efemeriden van de nacht}

Berekend vanuit de locatie geven zij zonsondergang, nautische schemering,
astronomische nacht en zonsopkomst --- en vooral hoeveel van de donkere nacht
de sessie werkelijk heeft bestreken. Een nacht die een uur te laat begint,
verliest een uur dat niet in te halen is, en dat leest nuttiger naast de
volgcijfers dan achteraf.

% ─────────────────────────────────────────────────────────────────
\chapter{De bestanden}
\label{ch:files}

\section{Wat een sessie schrijft}

\begin{center}
\rowcolors{2}{rowalt}{white}
\begin{tabularx}{\textwidth}{lX}
\toprule
\textbf{Extensie} & \textbf{Inhoud} \\
\midrule
\texttt{.dat} & De meetpunten: tijd, RK, DEC, standaardafwijkingen, status. \\
\texttt{.dti} & De synchronisatiemetingen tussen pc en montering. \\
\texttt{.fft} & De tijdens de sessie vastgelegde spectra. \\
\texttt{.env} & Temperatuur, luchtdruk, vochtigheid, als de montering ze
geeft. \\
\texttt{.log} & De gebeurtenissen: verbindingen, zwenkingen, meldingen,
onderbrekingen. \\
\bottomrule
\end{tabularx}
\end{center}

De kop van het \texttt{.dat}-bestand draagt de locatie en de tijdzone. Dat is
wat toelaat de schemeringen bij het terugkijken opnieuw te berekenen, lang
nadat de montering is losgekoppeld.

\section{Het nachtrapport}

Het wordt aan het eind van een sessie automatisch geschreven, naast het
\texttt{.dat}-bestand, als \textbf{één bestand per taal}. Er hoeft niets voor
gevraagd te worden.

\section{Terugkijken}

\textbf{Bestand $\rightarrow$ Logbestand openen} leest een afgeronde sessie en
bouwt de volledige analyse opnieuw op: statistieken per doel, aangestuurde
bewegingen, drift, spectra, efemeriden. De taaltabbladen tonen hetzelfde
rapport in alle drie de talen.

% ─────────────────────────────────────────────────────────────────
\chapter{Als er iets misgaat}
\label{ch:trouble}

\begin{center}
\rowcolors{2}{rowalt}{white}
\begin{tabularx}{\textwidth}{lX}
\toprule
\textbf{Verschijnsel} & \textbf{Waar te kijken} \\
\midrule
Geen netwerkverbinding & Controleer het adres en poort 3492, en of niet iets
anders de sessie al bezet houdt --- een 10Micron aanvaardt er maar één. \\
Seriële verbinding stil & Beide uiteinden moeten het over de snelheid eens
zijn. 10Micron-monteringen komen af fabriek op 9600 baud, en het menu van de
montering kan dat wijzigen. \\
Drift in declinatie & Voeg op een montering met model punten toe in de buurt
van het doel. De azimutschroef is niet het antwoord. \\
Matig oordeel, rustige montering & Lees hoofdstuk~\ref{ch:stats}: vaak is het
de lezing die fout is, niet de montering. \\
Onzinnige lopende STDEV & Normaal vlak na een zwenking of een
herpositionering: haar venster overspant de beweging. Het rapport laat die
waarden buiten beschouwing. \\
Geen efemeriden & De locatie is niet ingevuld. Voorkeuren, deel locatie. \\
\bottomrule
\end{tabularx}
\end{center}

% ─────────────────────────────────────────────────────────────────
\chapter{Naslag}
\label{ch:ref}

\section{Sneltoetsen}

\begin{center}
\rowcolors{2}{rowalt}{white}
\begin{tabularx}{\textwidth}{lX}
\toprule
\textbf{Toetsen} & \textbf{Handeling} \\
\midrule
\texttt{F1} & Hulp \\
\texttt{Ctrl+,} & Voorkeuren \\
\texttt{Ctrl+O} & Logbestand openen \\
\texttt{Ctrl+Q} & Afsluiten \\
\bottomrule
\end{tabularx}
\end{center}

\section{Taal}

Het menu \textbf{Taal} biedt Nederlands, Engels en Frans, plus automatische
herkenning vanuit het systeem. Elke regel staat in haar eigen taal geschreven:
wie in de verkeerde taal belandt, kan de huidige niet lezen om eruit te komen.
De wijziging gaat in bij de volgende start.

\section{Melden van problemen}

Loopt het programma vast, blokkeert het of wil het niet starten, dan kan het
dat zelf melden. De vraag wordt u eenmaal gesteld, en er wordt nooit iets
anders verstuurd: niet uw volggegevens, niet uw bestandsnamen, niet uw
gebruikersnaam. Alle paden worden geanonimiseerd.

\section{Versiegeschiedenis}

Het volledige overzicht van de wijzigingen, met de gemeten cijfers die ze
staven, staat in het bestand \texttt{CHANGELOG.md} dat met het programma wordt
meegeleverd.

\section{Herkomst}

MountMonitor is een zelfstandige herschrijving in Python en PyQt6 van
MountMonitor v3.37 van \textbf{Nicolàs de Hilster} (2018--2021), geschreven in
Java. Er staat geen regel code uit het oorspronkelijke programma in. De
MIT-licentie van dit programma dekt uitsluitend deze herschrijving; het
oorspronkelijke werk blijft onderworpen aan de voorwaarden van zijn auteur.
